\section{Method}
\label{sec:method}

\begin{figure*}[t]
\centering
\includegraphics[width=\textwidth]{foundation_ssc_instance_pipeline.png}
\caption{Foundation-model pipeline for instance-aware semantic scene completion (SSC). Monocular or multi-view OccuFly RGB observations are encoded by V-JEPA~2.1 or DINOv3 foundation features. Context and geometry adapters feed hybrid-view 3D lifting and axis-aware fusion, followed by spatial decoupling into semantic instance aggregation and 4D offset regression branches. Semantic class labels and instance offsets provide the respective supervision signals.}
\label{fig:architecture}
\end{figure*}

We evaluate frozen foundation-model representations on two aerial 3D tasks: metric depth estimation and semantic scene completion (SSC). We select V-JEPA~2.1 and DINOv3 as complementary video- and image-pretrained representations and organize the framework as a sequence of depth-readout and SSC-decoder comparisons. For MDE, we compare final-layer linear probes with multilevel DPT-style dense decoders, and examine camera Intrinsics-altitude-aware FiLM conditioning, video-based dense temporal attention (DTA), pose-aware multiview matching, and native MVSFormer++ geometry. For SSC, we retain CGFormer as a monocular control and evaluate FoundationSSC-based variants that combine frozen V-JEPA~2.1 or DINOv3 context with FiLM-DPT monocular depth, calibrated MVSFormer++ geometry, or privileged GT depth, followed by either the retained FSSC voxel head or VoxDet decoding. Together with the MDE variants above, these configurations examine frozen-representation transfer, dense and metadata -conditioned depth decoding, temporal aggregation, geometry source, and voxel decoding. Concretely we seperate four questions: whether geometry is accessible from a frozen representation (RQ1), whether metadata helps beyond decoder capacity (RQ2), whether calibrated motion compensation improves temporal geometry (RQ3), and whether SSC is limited by depth, lifting, or semantic taxonomy (RQ4).

\subsection{Tasks and Information Pathways}
Each target observation has image $I_t$, intrinsics $K_t$, altitude $h_t$, and pose $T_t^{w\rightarrow c}$. Supervised training additionally uses depth $D_t$, its validity mask $M_t$, and semantic voxel labels $Y_t$. MDE predicts camera-axis depth $\hat D_t$ in meters. SSC predicts a categorical label $\hat Y_t(v)$ for each voxel $v$, with label zero representing empty space and labels 1--21 representing occupied semantic categories. Ground-truth (GT) depth and semantic labels supply supervision.

The pipeline separates context, geometry, and decoding. A frozen encoder supplies visual features. A context head maps those features to a semantic representation, while a geometry head or frozen MVS provider supplies depth information. A hybrid view transformer lifts context into a common voxel volume and fuses the resulting features. A voxel decoder then predicts occupancy and semantics. Figure~\ref{fig:architecture} shows these boundaries for depth transfer and the final multiview treatment. The division makes it possible to describe what changes across configurations without suggesting that every comparison isolates one factor.

For a target with neighboring views, define
\begin{equation}
\mathcal V_t=\{(I_i,K_i,T_i^{w\rightarrow c})\}_{i\in\{t\}\cup\mathcal S_t},
\end{equation}
where $\mathcal S_t$ contains source views. The monocular path uses only the target image and the metadata supplied to its head. The final multiview path uses one target and four calibrated sources. The temporal depth diagnostic uses a shorter target-centered neighborhood; it is a separate experiment and does not define the final SSC provider.

\subsection{Frozen Image Features}
The V-JEPA~2.1 ViT-Gigantic image encoder takes $384\times384$ inputs with patch size 16. It exposes $24\times24$ maps with 1664 channels. The four-level decoder reads zero-indexed blocks $\{11,23,37,47\}$; the last-layer probe reads block 47 alone. DINOv3 ViT-H+/16 uses the same target resolution and patch grid, with 1280-channel maps from blocks $\{7,15,23,31\}$.

Let $F_t^{(l)}=E^{(l)}(I_t)$ denote the selected frozen map. The foundation encoder is frozen: its features participate in the forward pass, but its parameters are not updated by gradients from the downstream task loss.

\begin{figure}[t]
\centering
\includegraphics[width=\columnwidth]{vjepa_dinov3_feature_diagnostics.png}
\caption{Matched feature and depth diagnostics for V-JEPA~2.1 and DINOv3 on OccuFly scene~08 at 50~m (frame 000091). The final frozen maps are block~47 for V-JEPA and block~31 for DINOv3. DPT level~0 is the higher-resolution branch derived from the earliest selected backbone level and is architecturally positioned to retain finer spatial variation. Level~3 is the coarsest branch derived from the deepest selected level and has the largest effective context. These are architectural roles, not semantic interpretations of PCA colors. Each PCA is fitted independently and supports only within-panel inspection; colors are not comparable across models or levels. Feature RMS shows within-backbone activation magnitude using independently normalized low-to-high scales. }
\label{fig:backbone_feature_diagnostics}
\end{figure}

\subsection{Depth Heads and Metric Readout}
\subsubsection{Linear and DPT decoding}
The linear probe applies batch normalization and a $1\times1$ prediction layer to the final feature map. The head produces 256 depth-channel responses and a dense metric readout after spatial resizing. Its small capacity provides a baseline for geometry that can be extracted with limited adaptation.

The four-level DPT head projects selected maps into a common decoder representation and reassembles them at different spatial scales. Convolutional refinement progressively combines coarse and fine maps before depth prediction \cite{ranftl2021dpt}. Although all selected transformer maps have the same patch-grid size, their encoder depth differs; reassembly supplies the spatial pyramid needed by the dense decoder. The head therefore changes both the feature levels available and the reconstruction capacity relative to a final-layer linear probe.

\subsubsection{Depth-distribution readout}
At each pixel $u$, an image-depth head predicts $K=256$ responses
$z_{t,k}(u)$ associated with fixed metric-depth values $d_k$. A
head-specific nonnegative mapping followed by normalization produces
depth weights
\begin{equation}
p_{t,k}(u)=
\frac{a(z_{t,k}(u))}
     {\sum_j a(z_{t,j}(u))},
\qquad
\hat D_t(u)=\sum_k p_{t,k}(u)d_k .
\label{eq:depth}
\end{equation}
The reported image-depth map is therefore the expectation of the
discrete metric-depth representation. MVSFormer++ instead
returns an explicit categorical posterior over its native depth
hypotheses; for SSC, this posterior is projected onto the lifting grid
as described below.

\subsubsection{Camera-conditioned FiLM}
The metadata vector uses altitude and resized intrinsics:
\begin{equation}
m_t=\left[\frac{h_t-40}{10},\frac{f_x}{W},\frac{f_y}{H},
\frac{c_x}{W}-\frac12,\frac{c_y}{H}-\frac12\right].
\label{eq:metadata}
\end{equation}
A multilayer perceptron embeds $m_t$ into 128 dimensions. At each DPT level, channel-wise projections produce scale and shift vectors, giving
\begin{equation}
\widetilde R_l=R_l\odot(1+\gamma_l(m_t))+\beta_l(m_t).
\label{eq:film}
\end{equation}
The affine projections start near zero so the initial path remains close to the unconditioned decoder. Metadata thus modulates the interpretation of visual features while leaving the frozen encoder untouched. Camera conditioning enters only methods named FiLM-DPT.

\subsection{Temporal Aggregation and Pose-Aware Diagnostic}
\begin{figure*}[!t]
\centering
\includegraphics[width=\textwidth]{dtavideo/dta_raft_exact.pdf}
\caption{Flow-aligned DTA diagnostic for the four-frame V-JEPA~2.1 clip. The raw tubelet PCA maps precede alignment; feature change is their channel-wise $\ell_2$ difference, normalized at the 99th percentile. RAFT flow and validity backward-warp the past and future tubelets into target coordinates before DTA, whose blue/red preference favors tubelet 0/1. Raw, aligned, and DTA maps share one PCA basis; the predicted depth uses a 0--80~m scale.}
\label{fig:dta_raft_compact}
\end{figure*}

Video mode uses a four-frame clip with offsets $[-1,0,0,+1]$. The duplicated target yields an even length compatible with tubelet size two. Frame offsets are defined by sequence index, because measured capture timestamps are not available in the audited preprocessing records. The frozen video encoder returns two temporal tubelet positions at each spatial patch location.

A video linear adapter folds the tubelet dimension into channels. Dense Temporal Attention (DTA) instead learns a temporal readout separately at each spatial cell and each selected encoder level. With temporal tokens $G_{l,\tau}(u)$, the adapter has the form
\begin{equation}
\bar F_l(u)=\mu_l(u)+\alpha_l\bigl(A_l(G_l(u))-\mu_l(u)\bigr),
\label{eq:dta}
\end{equation}
where $\mu_l$ is temporal mean pooling, $A_l$ is learned-query attention, and the residual gate starts near mean pooling. The attention module uses a 512-dimensional adapter with eight heads. Its output has the same two-dimensional spatial form as an image feature map and can feed FiLM-DPT.

The implemented flow-aligned diagnostic addresses the same-cell assumption explicitly. Frozen RAFT \cite{teed2020raft} predicts target-to-frame flow $F_{t\rightarrow j}$ and reverse flow for every RGB frame. Out-of-view samples and forward--backward-inconsistent matches are invalidated \cite{meister2018unflow}. Because one V-JEPA token spans two frames, the two valid frame flows are averaged to obtain tubelet flow $\bar F_\tau$. For target coordinate $u$, source-tubelet features are then sampled as
\begin{equation}
\widetilde G_{l,\tau}(u)=G_{l,\tau}\!\left(u+\bar F_\tau(u)\right),
\label{eq:raft_warp}
\end{equation}
using bilinear backward sampling with \texttt{align\_corners=False}. The aligned DTA masks invalid tubelets in attention and replaces $\mu_l$ in Eq.~\eqref{eq:dta} by a validity-weighted mean. The aligned Linear Probe Video DTA control (MDE-2 in Table~\ref{tab:mde}) instead zeroes invalid samples before temporal concatenation, preserving a linear depth readout. RAFT, V-JEPA~2.1, and the warping operation are frozen or deterministic; only the downstream temporal adapter and depth head are trained.

Local temporal attention does not align features when the camera moves, because the same patch coordinate can correspond to different physical surfaces. The pose-aware component of Video DTA-FiLM-MVS (MDE-3 in Table~\ref{tab:mde}) therefore uses calibrated, depth-conditioned projective warping.

For homogeneous target pixel $\tilde u$ and depth hypothesis $d_k$, the target-camera point and its projection into source view $i$ are
\begin{equation}
x_{t,k}=d_kK_t^{-1}\tilde u,\qquad
q_{i,k}=\pi\!\left(K_i(R_{it}x_{t,k}+t_{it})\right),
\label{eq:projection}
\end{equation}
where
$T_{t\rightarrow i}
=T_i^{w\rightarrow c}(T_t^{w\rightarrow c})^{-1}
=[R_{it}\mid t_{it}]$
is the target-to-source camera transform and $\pi$ performs perspective division.

The source feature map is then warped into the target-view depth volume by bilinear sampling:
\begin{equation}
\widetilde F_{i,k}(u)
=
\operatorname{bilinear}\!\left(F_i,q_{i,k}\right).
\label{eq:mvs_warp}
\end{equation}
The warped source features are compared with the target features and combined across two source views to form a depth-indexed matching volume. A three-dimensional convolutional head regularizes this volume, and FiLM conditions the resulting depth responses. Unlike the preceding RAFT warp, this projective warp is determined by camera calibration, relative pose, and the tested depth hypothesis.

\subsection{Frozen Native Multiview Geometry}
The final multiview treatment (SSC-5 in Table~\ref{tab:ssc}) uses a frozen DINOv3-backed MVSFormer++ provider with one target and four calibrated source images. MVSFormer++ operates on $768\times1024$ images, while its foundation-feature branch uses $336\times448$ inputs. The retained geometry checkpoint is loaded strictly, and neither the DINOv3 backbone nor MVSFormer++ is updated during SSC training.

MVSFormer++ returns depth hypotheses $z_j(u)$ and their categorical probabilities $p_j(u)$. For our OccuFly SSC configuration, the LSS module discretizes camera-axis depth over $[0,64]$~m into $K=256$ bins of width 0.25~m, with centers $b_k$ from 0.125 to 63.875~m. This ray-depth discretization is distinct from the 0.5-m resolution of the OccuFly occupancy grid. We therefore linearly redistribute each MVS probability between its two neighboring lifting centers:
\begin{equation}
\widetilde P_k(u)=\sum_j p_j(u)w_k\!\left(z_j(u)\right),
\qquad
P_k(u)=\frac{\widetilde P_k(u)}
{\sum_l\widetilde P_l(u)} ,
\label{eq:posterior}
\end{equation}
where $w_k(z)$ is the linear interpolation weight for center $b_k$. Hypotheses outside the lifting range receive zero weight. If no probability remains within the range, the ray is marked invalid and does not contribute to lifting.

The context image is derived from the same target crop as the MVS input, with its intrinsics scaled to the corresponding resolution. This keeps the context features and projected depth posterior spatially aligned. Expected depth and photometric confidence are retained only as diagnostics; SSC lifting uses the complete projected distribution $P_k(u)$.

\subsection{Probability-Aware Lifting and Voxel Fusion}
The semantic context head consumes the target's frozen image features and produces a trainable DPT context map $C_t(u)$. For each image coordinate and lifting depth center, a camera ray maps context to a three-dimensional location. Schematically, the lift--splat contribution to voxel $v$ is
\begin{equation}
V(v)=\sum_{u,k}\mathbf1[\nu(x_{t,k}(u))=v]P_k(u)C_t(u),
\label{eq:lift}
\end{equation}
where $x_{t,k}(u)=b_kK_t^{-1}u$ and $\nu$ assigns a position to the configured voxel grid. Equation~\eqref{eq:lift} describes probability-weighted accumulation; the implemented hybrid view transformer additionally combines this dense pathway with proposal-guided voxel features and axis-aware fusion from the retained FoundationSSC modules.

The fused representation has 128 channels at $96\times64\times64$, corresponding to a reduced grid before full-resolution prediction. In the final multiview treatment, the retained view transformer and fusion pathway are frozen. Context and VoxDet decoding are trained while geometry and the lifting modules remain fixed. Context trains through the dense lift--splat and fusion path. Inputs to the frozen sparse voxel-attention branch are detached, preserving its forward features without backpropagating through that branch.

\subsection{VoxDet Decoding and VoxNT Targets}
VoxDet processes the fused volume with a shared three-dimensional encoder and task-specific branches. One branch predicts semantic features and logits, while another predicts six nonnegative directional offsets. Predicted offsets guide sampling and aggregation of semantic features along the voxel axes before final classification \cite{li2025voxdet}. The scored output remains a semantic occupancy grid. Directional predictions do not establish instance detection accuracy, and no instance average precision is claimed.

VoxNT generates supervision from semantic labels without learning additional parameters. For a voxel coordinate $i$ along an axis, the inclusive positive-direction run length is
\begin{equation}
r_+(i)=1+\mathbf1[Y(i)=Y(i+1)]r_+(i+1),
\label{eq:run}
\end{equation}
with terminal run length one. Applying the recurrence in both directions along three axes yields six targets. Each is normalized by its corresponding full-grid axis length. This is an axial same-label run, not Euclidean distance to the nearest object boundary or a connected-component box. A gap or class change ends the run, while touching same-label objects can remain indistinguishable.

Our implemented supervision enlarges predicted reduced-grid offsets to the full target grid and uses class-balanced Smooth-L1 loss with transition parameter 0.1. Eligible categories are tree, ground obstacle, person, bicycle, vehicle, building, cable, construction, crane, and truck. Losses are averaged within each present eligible class and then across those classes, reducing domination by large regions. All categories still receive semantic occupancy supervision. Target normalization uses the full-grid axis lengths before any resolution change.

Feature sampling uses the center-aligned coordinate
\begin{equation}
g(i;L)=2\frac{i+1/2}{L}-1
\label{eq:sample}
\end{equation}
for an axis of length $L$. This keeps voxel-center indexing consistent with the sampler's alignment convention. Tensor storage order is respected when coordinates are passed to the three-dimensional sampler. The treatment combines correct target scaling, consistent sampling, full-resolution supervision, and a long-tail objective; the available completed comparisons cannot assign its overall gain to one of these corrections alone.

Ground-truth VoxNT targets are used only during training. At inference, offsets are predicted from features, and no semantic target tensor is supplied. The phrase instance-aware describes the intended structural aggregation, with the semantic-label limitation above; it should not be read as verified separation of every physical instance.

\subsection{Occupancy and Long-Tail Semantic Objective}
Let $\ell_c(v)$ be the 22-class logit vector. Binary occupancy is trained from the stable occupied-versus-empty logit
\begin{equation}
\ell_{\rm occ}(v)=\log\sum_{c=1}^{21}\exp\ell_c(v)-\ell_0(v).
\end{equation}
A weighted binary cross-entropy (BCE) term supervises every valid voxel. Its positive weight is $\min(25,n_0/\sum_{c>0}n_c)$ from training counts. Conditional semantic focal cross-entropy is computed only at ground-truth occupied voxels, with focal exponent two. Semantic weights use square-root inverse training frequency, clipped to $[0.25,8]$ and then normalized to mean one over represented occupied classes. Empty voxels consequently supervise occupancy without overwhelming conditional category recognition.

A macro soft Dice term averages over semantic classes present in each batch. The main objective is
\begin{equation}
\mathcal L_{\rm main}=\mathcal L_{\rm occ}
+\mathcal L_{\rm focal}+0.5\mathcal L_{\rm Dice}.
\label{eq:mainloss}
\end{equation}
The auxiliary head uses corresponding weights $0.25$, $0.25$, and $0.10$, and full-resolution VoxNT has weight one. These weights are implementation settings rather than conclusions of a hyperparameter ablation.

The monocular VoxDet treatment additionally uses a metric-depth teacher term with weight 0.1 and a weak scheduled SSC-to-depth gradient. That gradient is zero in epoch zero, ramps over three epochs, and is capped at 0.1. Its purpose is to let the voxel decoder stabilize before limited adaptation of the pretrained depth head. The frozen multiview treatment has no depth loss and no such gradient schedule: its provider is run without gradients. All semantic, geometric, and offset masks exclude ignored labels. The distinction between these gradient paths is part of the method definition and limits direct architectural attribution from the final scores.
